\heading{量子力学A-23秋-期末}
\noteinfo{整理自 Quantum Mechanics: Fall 2023 Final Exam: Brief Solutions；题面和解答按项目格式重写。}

\begin{question}
粒子在周长 $L$ 的圆环上运动，加入 $V(x)=V_0\cos(2\pi x/L)$。
\begin{enumerate}
  \item 写出自由粒子本征值和本征函数。
  \item 求能量修正，特别讨论偶、奇宇称子空间。
  \item 对变分波函数 $\psi=A+B\cos(2\pi x/L)$ 计算最小能量。
\end{enumerate}
\begin{solution}
自由本征态
\[
  \psi_n={1\over\sqrt L}\ee^{\ii2\pi nx/L},\qquad
  E_n^{(0)}={\hbar^2\over2m}\left({2\pi n\over L}\right)^2.
\]
在实基底中
\[
  \psi_{0,e}=L^{-1/2},\quad
  \psi_{n,e}=\sqrt{2/L}\cos{2\pi nx\over L},\quad
  \psi_{n,o}=\sqrt{2/L}\sin{2\pi nx\over L}.
\]
势为偶函数，偶、奇子空间不混合。基态二阶修正
\[
  E_0=-{V_0^2mL^2\over4\pi^2\hbar^2}+O(V_0^3).
\]
偶态 $n=1$ 有一阶修正 $V_0/2$，二阶修正为
\[
  \Delta E_{1,e}^{(2)}={5V_0^2mL^2\over24\pi^2\hbar^2};
\]
奇态 $n=1$ 一阶为零，二阶修正为
\[
  \Delta E_{1,o}^{(2)}=-{V_0^2mL^2\over24\pi^2\hbar^2}.
\]
一般 $n>1$ 只耦合 $n\pm1$，二阶项为
\[
  \Delta E_n^{(2)}={V_0^2mL^2\over4\pi^2\hbar^2}{1\over4n^2-1}.
\]
变分态归一化后
\[
  E(A,B)={E_1^{(0)}B^2/2+V_0AB\over A^2+ B^2/2}.
\]
令 $r=B/A$ 为变分参数，最小化
\[
  E(r)={E_1^{(0)}r^2/2+V_0r\over1+r^2/2}
\]
得到
\[
  E_{\min}={1\over2}\left(E_1^{(0)}-\sqrt{(E_1^{(0)})^2+2V_0^2}\right).
\]
\end{solution}
\end{question}

\begin{question}
二维谐振子受旋转外力
\[
  V(t)=-f[x\sin\Omega t+y\cos\Omega t].
\]
\begin{enumerate}
  \item 写出从基态到第一激发态的跃迁振幅。
  \item 求从基态到第一激发能级的总跃迁概率。
  \item 求到 $\ket{1,1}$ 的最低非零阶跃迁概率。
  \item 定义旋转坐标，写出旋转参考系中的哈密顿量。
\end{enumerate}
\begin{solution}
使用升降算符
\[
  x=\sqrt{\hbar\over2m\omega}(a_x+a_x^\dagger),\qquad
  y=\sqrt{\hbar\over2m\omega}(a_y+a_y^\dagger).
\]
一阶振幅为
\[
  c_{10}^{(1)}={\ii f\over\hbar}\sqrt{\hbar\over2m\omega}
  \int_0^t\ee^{\ii\omega t'}\sin\Omega t'\,\dd t',
\]
\[
  c_{01}^{(1)}={\ii f\over\hbar}\sqrt{\hbar\over2m\omega}
  \int_0^t\ee^{\ii\omega t'}\cos\Omega t'\,\dd t' .
\]
总跃迁概率为 $|c_{10}^{(1)}|^2+|c_{01}^{(1)}|^2$。到 $\ket{1,1}$ 需二阶过程，振幅等于 $x$、$y$ 两个一次跃迁振幅的有序积分和。

旋转坐标
\[
  \tilde x=x\cos\Omega t-y\sin\Omega t,\qquad
  \tilde y=x\sin\Omega t+y\cos\Omega t
\]
中哈密顿量变为
\[
  \tilde H=H_0-\Omega L_z-f\,\tilde y
\]
（取决于旋转方向约定，$L_z$ 项符号相应改变）。这给出了不含时问题。
\end{solution}
\end{question}

\begin{question}
两个全同粒子在一维圆环上运动，
\[
  H_0={p_1^2+p_2^2\over2m}.
\]
\begin{enumerate}
  \item 写出无自旋玻色子的基态、第一激发态、第二激发态。
  \item 写出无自旋费米子的基态、第一激发态、第二激发态。
  \item 写出自旋 $1/2$ 玻色子、费米子的基态和第一激发态。
  \item 在第一激发态中计算 $\cos[2\pi(x_1-x_2)/L]$ 的期望值。
  \item 在基态中测量 $S_{1x}+S_{2x}$，求可能结果和概率。
  \item 加入 $H_1=\lambda\delta(x_1-x_2)$，求自旋 $1/2$ 玻色子在前述能级中的最低非零阶能量修正。
  \item 推导两个自旋 $1/2$ 玻色子基态能量的精确方程。
\end{enumerate}
\begin{solution}
单粒子本征态为
\[
  \varphi_k(x)={1\over\sqrt L}\ee^{\ii2\pi kx/L},\qquad
  \epsilon_k={\hbar^2\over2m}\left({2\pi k\over L}\right)^2 .
\]
无自旋玻色子基态为 $\varphi_0(1)\varphi_0(2)$，能量 $0$。第一激发态为
\[
  {1\over\sqrt2}\bigl[\varphi_{\pm1}(1)\varphi_0(2)
  +\varphi_0(1)\varphi_{\pm1}(2)\bigr],
  \qquad E={2\pi^2\hbar^2\over mL^2},
\]
第二激发态可取 $\varphi_{\pm1}(1)\varphi_{\pm1}(2)$ 和
$(\varphi_1(1)\varphi_{-1}(2)+\varphi_{-1}(1)\varphi_1(2))/\sqrt2$，
能量为两倍第一激发能。无自旋费米子基态必须占据两个不同动量态：
\[
  {1\over\sqrt2}\bigl[\varphi_0(1)\varphi_{\pm1}(2)
  -\varphi_{\pm1}(1)\varphi_0(2)\bigr],
  \qquad E={2\pi^2\hbar^2\over mL^2},
\]
第一激发态为
$(\varphi_1(1)\varphi_{-1}(2)-\varphi_{-1}(1)\varphi_1(2))/\sqrt2$，
能量为其两倍。

自旋 $1/2$ 时，玻色子要求总波函数对称，故空间对称态配自旋三重态，
空间反对称态配自旋单态；费米子相反。于是玻色子基态为
$\varphi_0\varphi_0\ket{1,m}$，$m=1,0,-1$；费米子基态为
$\varphi_0\varphi_0\ket{0,0}$。第一激发能级中同时有空间对称和空间反对称
轨道态，分别按上述规则配自旋态。

在第一激发态中，
\[
  \expval{\cos {2\pi(x_1-x_2)\over L}}=
  \begin{cases}
  1/2,&\text{空间波函数对称},\\
  -1/2,&\text{空间波函数反对称}.
  \end{cases}
\]
基态自旋测量 $S_{1x}+S_{2x}$ 的概率只由自旋态决定。自旋单态只能得到
$0$，概率为 $1$；三重态 $\ket{1,m_z}$ 中，$\ket{1,\pm1_z}$ 给出
$+\hbar,0,-\hbar$ 的概率 $1/4,1/2,1/4$，而 $\ket{1,0_z}$ 给出
$+\hbar,-\hbar$ 的概率各 $1/2$。

加入 $H_1=\lambda\delta(x_1-x_2)$ 时，接触相互作用只作用在空间对称部分。
玻色子基态的一阶修正为
\[
  \Delta E_0^{(1)}={\lambda\over L}.
\]
第一激发能级中，空间对称轨道态的一阶修正为 $2\lambda/L$，空间反对称轨道态
因在 $x_1=x_2$ 处为零而一阶修正为 $0$。

对两个自旋 $1/2$ 玻色子的精确基态，可取质心 $X=(x_1+x_2)/2$ 与相对坐标
$x=x_2-x_1$。基态的质心动量为零；相对运动满足
\[
  \left(-{\hbar^2\over m}{\dd^2\over\dd x^2}+\lambda\delta(x)\right)\phi(x)=E\phi(x),
  \qquad \phi(x)=\phi(-x)=\phi(x+L).
\]
若 $E=\hbar^2k^2/m$，在 $0<x<L/2$ 可写
$\phi=A\cos[k(x-L/2)]$，跳变条件给
\[
  {2\hbar^2k\over m\lambda}=\cot{kL\over2}.
\]
这就是基态能量的隐式方程；$\lambda<0$ 时令 $k=\ii\kappa$ 得到束缚态形式。
\end{solution}
\end{question}
